\subsection{用光子系统检验不变质量}
单个真空光子的静质量为零，并不意味着多个光子组成的系统不变质量也为零。两光子的能量分别为 $E_1,E_2$，传播方向夹角为 $\theta$，总动量模平方为
\[
 P^2c^2=E_1^2+E_2^2+2E_1E_2\cos\theta.
\]
由总能量 $E=E_1+E_2$，得到
\[
 M^2c^4=E^2-P^2c^2=2E_1E_2(1-\cos\theta).
\]
同向时 $M=0$，不存在总动量为零的惯性系；反向且等能 $E_1=E_2=\varepsilon$ 时，总动量为零，$M=2\varepsilon/c^2$。这里的质量属于系统，不是给每个光子分配了静质量。计算碰撞阈值时使用总四动量，正是为了保留这类内部相对运动的能量。
